\begin{figure}[htbp]
\centering
\begin{tikzpicture}[>=Stealth,state/.style={circle,draw,minimum size=7mm,inner sep=1pt},every node/.style={font=\small}]
\node at (1.9,1.15) {Active slot};
\node at (7.25,1.15) {Unused slot};
\draw[dashed,gray] (1.9,-0.9)--(1.9,0.75);
\draw[dashed,gray] (7.25,-0.9)--(7.25,0.75);
\node[state] (activefrom) at (0.5,0) {$p$};
\node[state] (activeto) at (3.3,0) {$q$};
\draw[->,blue!65!black] ([yshift=2pt]activefrom.east)--([yshift=2pt]activeto.west);
\draw[<-,blue!65!black] ([yshift=-2pt]activefrom.east)--([yshift=-2pt]activeto.west);
\node[state] (unusedfrom) at (5.85,-0.6) {$p$};
\node[draw,rectangle,minimum size=2mm,inner sep=0pt] (leaf) at (5.85,0.3) {};
\node[above=2pt of leaf] {fresh leaf};
\node[state] (unusedto) at (8.65,0.3) {$q$};
\draw[->,blue!65!black] ([yshift=2pt]leaf.east)--([yshift=2pt]unusedto.west);
\draw[<-,blue!65!black] ([yshift=-2pt]leaf.east)--([yshift=-2pt]unusedto.west);
\node at (0.6,-1.25) {source cell};
\node at (3.25,-1.25) {target cell};
\node at (5.95,-1.25) {source cell};
\node at (8.55,-1.25) {target cell};
\end{tikzpicture}
\caption{One candidate transition slot with $d=\rightarrow$, with its two travel
lanes. The target end is always attached to state $q$. Only the source
attachment changes: an actual transition uses state $p$, and an unused
slot uses its own new leaf, adding no state-to-state connection.
Thus the target cell needs no information
about the source symbol. Other slots and local cyclic joins are omitted;
lane arrows indicate traversal, not two machine transitions.}
\label{fig:candidate-slot}
\end{figure}
